\documentclass[11pt,a4paper]{article}
\usepackage{fontspec,amsmath,amssymb,amsthm,microtype}
\setmainfont{lmroman10-regular.otf}[BoldFont=lmroman10-bold.otf,ItalicFont=lmroman10-italic.otf,BoldItalicFont=lmroman10-bolditalic.otf]
\setmonofont{lmmono10-regular.otf}[Scale=MatchLowercase]
\usepackage[margin=25mm]{geometry}
\usepackage[unicode,hidelinks]{hyperref}
\setlength{\emergencystretch}{3em}
\setlength{\parskip}{0.4em}
\setlength{\parindent}{0pt}
\allowdisplaybreaks[2]
\raggedbottom
\title{\texorpdfstring{Relative entropy along an identified density path}{Relative entropy along an identified density path}}
\author{LangevinLean mathematical exposition}
\date{}
\begin{document}\maketitle
\noindent Source module: \texttt{Unadjusted\allowbreak{}Langevin/\allowbreak{}LSI/\allowbreak{}Entropy\allowbreak{}Flow.\allowbreak{}lean}. The accompanying HTML pairs this mathematical exposition with the exact Lean statements and records the scope of the formalization.


Density entropy is the real conversion of the actual KL divergence of the tilted measure. With finite-entropy guards, LSI and a supplied dissipation inequality give exponential entropy decay along this path.


\section{\texorpdfstring{Define actual density entropy with its finiteness guard}{Define actual density entropy with its finiteness guard}}


Define the real entropy of a specified relative density as \(H_\nu(r)=\operatorname{KL}(r\nu\mid\nu)_{\mathbb R}\). It uses the same \texttt{tilted\allowbreak{}Measure} and actual Mathlib KL as the LSI definition. Applications must separately establish finite KL before using it as an ordinary entropy value.



{\raggedright\noindent Lean declarations: \texttt{density\allowbreak{}Entropy}.\par}

\section{\texorpdfstring{Apply the finite real LSI at every time}{Apply the finite real LSI at every time}}


At each interpolation time, positivity, \(C^1\) regularity, normalization, and finite KL/\allowbreak{}Fisher allow \texttt{Is\allowbreak{}Log\allowbreak{}Sobolev.\allowbreak{}entropy\_\allowbreak{}le\_\allowbreak{}fisher} to give \(2\alpha H(t)\le I(t)\). Combine this with the explicit right-dissipation identity and entropy continuity in \texttt{entropy\_\allowbreak{}decay\_\allowbreak{}of\_\allowbreak{}dissipation}. The conclusion is decay at rate \(2\alpha\) for that admissible density path; identifying an actual evolution with the path remains an application obligation.



{\raggedright\noindent Lean declarations: \texttt{relative\_\allowbreak{}entropy\_\allowbreak{}decay\_\allowbreak{}of\_\allowbreak{}dissipation}.\par}

\end{document}
